\section{Related Work}
\label{sec:related_work}

\para{Binary detection and source attribution.} Surveys of machine-generated text detection organize the space into statistical, neural, watermarking, and human-assisted methods, and emphasize that robustness remains unresolved \cite{wu2023survey}. \semeval established the strongest benchmark in our review for moving beyond binary labeling, with subtasks for human-vs-machine detection, exact generator attribution, and human-machine boundary detection \cite{wang2024semeval}. Our work builds directly on that framing, but we ask a different question: instead of starting from generator identity, we start from broad production modes and ask whether an LLM reader can recover them under approximate content control.

\para{Fine-grained process labels.} \llmdetectaive shows that machine-generated, machine-humanized, and human-machine-polished texts can be separated in a four-way setup \cite{abassy2024llmdetectaive}. Ji et al.\ argue that ambiguity itself deserves explicit treatment, since many texts do not fit cleanly into binary human-vs-AI labels \cite{ji2024notjust}. \textsc{TextMachina} complements this line by providing infrastructure for generating attribution and mixed-authorship datasets when public benchmarks are insufficient \cite{sarvazyan2024textmachina}. We share the process-label motivation of these papers, but unlike them we keep content anchors roughly matched across multiple modes and pair end-task classification with embedding-space tests.

\para{Spoken and behavioral signals.} Not all production traces are tied to final-text authorship. \swab frames spoken-to-written conversion as a document-level transformation problem and shows that transcript artifacts form a coherent mode that LLMs can rewrite \cite{liu2024recording}. Kundu et al.\ show that AI assistance can also be detected from keystroke dynamics rather than final text alone \cite{kundu2024keystroke}. These papers motivate two of our mode families: dictated/spoken rewrites and keyboard-noisy text. Our study focuses on text-only evidence, so it should be read as complementary to richer behavioral detection setups rather than a replacement for them.

\para{Human- and model-authored comparison corpora.} \hcthree provides a practical source of matched human and ChatGPT answers across multiple QA domains \cite{guo2023hc3}. We use \hcthree as the anchor corpus for our controlled benchmark because it gives us real human answers that can be transformed into multiple process variants. This lets us vary production mode while holding topic and intent approximately fixed, which is the key design element missing from prior work.
